% Part: normal-modal-logic
% Chapter: sequent-calculus
% Section: rules-for-K

\documentclass[../../../include/open-logic-section]{subfiles}

\begin{document}

\olfileid{nml}{seq}{rul}

\olsection{\Log{K} కోసం నియమాలు}

సాధారణ ప్రతిజ్ఞావాక్య సంధాయకాల నియమాలు సాధారణ సీక్వెంట్
కలనశాస్త్రం~$\Log{LK}$లోని నియమాలే. స్వీకృతాలూ అవే:
$!A \Sequent !A$ రూపంలోని ఏ సీక్వెంట్ అయినా స్వీకృతం.

మోడల్ కారకం\iftag{prvBox}{\iftag{prvDiamond}{లైన~$\Box$
మరియు}{మైన~$\Box$}}{}\iftag{prvDiamond}{~$\Diamond$}{} కోసం
కింది అదనపు \iftag{notprvBox,notprvDiamond}{నియమం}{నియమాలు}
ఉన్నాయి:
\[
  \iftag{prvBox}{\iftag{prvDiamond}{% <> and [] primitive
    \Axiom$\Gamma \fCenter \Delta, !A$
    \RightLabel{$\Box$}
    \UnaryInf$\Box\Gamma \fCenter \Diamond\Delta, \Box!A$
    \DisplayProof
    \qquad
    \Axiom$!A, \Gamma \fCenter \Delta$
    \RightLabel{$\Diamond$}
    \UnaryInf$\Diamond!A, \Box\Gamma \fCenter \Diamond\Delta$
    \DisplayProof
}{% only Box primitive
\Axiom$\Gamma \fCenter !A$
\RightLabel{$\Box$}
\UnaryInf$\Box\Gamma \fCenter \Box!A$
\DisplayProof
}}{% only <> primitive
  \Axiom$!A \fCenter \Delta$
  \RightLabel{$\Diamond$}
  \UnaryInf$\Diamond!A \fCenter \Diamond\Delta$
  \DisplayProof
}
\]
ఇక్కడ, \iftag{prvBox}{$\Gamma$ నుంచి, ఆ $\Gamma$లోని ప్రతి
\tetoken{సూత్రం}{!!{formula}} ముందు $\Box$ ఉంచితే
వచ్చే \tetoken{సూత్రాల}{!!{formula}s} అనుక్రమమే
$\Box\Gamma$\iftag{prvDiamond}{; అలాగే
}{}}{}\iftag{prvDiamond}{$\Delta$ నుంచి, ఆ $\Delta$లోని ప్రతి
\tetoken{సూత్రం}{!!{formula}} ముందు $\Diamond$ ఉంచితే
వచ్చే \tetoken{సూత్రాల}{!!{formula}s} అనుక్రమమే
$\Diamond\Delta$}{}. \iftag{notprvDiamond}{$\Box$ నియమపు
పూర్వపక్షం కుడివైపున గరిష్ఠంగా ఒక్క
\tetoken{సూత్రమే}{!!{formula}}~$!A$ ఉండాలి. }{}%
\iftag{notprvBox}{$\Diamond$ నియమపు పూర్వపక్షం
ఎడమవైపున గరిష్ఠంగా ఒక్క
\tetoken{సూత్రమే}{!!{formula}}~$!A$ ఉండాలి. }{}%
\iftag{prvBox}{$\Gamma$\iftag{prvDiamond}{~మరియు~}{}}{}\iftag{prvDiamond}{$\Delta$}{}
ఖాళీగా ఉండవచ్చు; అప్పుడు నిష్కర్ష సీక్వెంట్‌లోని
సంబంధిత భాగం
\iftag{prvBox}{$\Box\Gamma$\iftag{prvDiamond}{~మరియు~}{}}{}\iftag{prvDiamond}{$\Diamond\Delta$}{}
కూడా ఖాళీగా ఉంటుంది.

\iftag{prvBox}{కుడివైపున $\Box$ను\iftag{prvDiamond}{,
అలాగే }{}}{}\iftag{prvDiamond}{ఎడమవైపున $\Diamond$ను}{}
ఒక్క \tetoken{సూత్రం}{!!{formula}}~$!A$కే చేర్చాలనే
పరిమితి అవసరం. ఒకవేళ
\iftag{prvBox}{కుడివైపున ఎన్ని
\tetoken{సూత్రాలకైనా}{!!{formula}s} $\Box$ను
చేర్చనిస్తే\iftag{prvDiamond}{ లేదా }{}}{}\iftag{prvDiamond}{ఎడమవైపున ఎన్ని
\tetoken{సూత్రాలకైనా}{!!{formula}s} $\Diamond$ను
చేర్చనిస్తే}{} కింది వాటిని
\tetoken{వ్యుత్పాదించగలం}{!!{derive}}:
  \[
    \iftag{prvBox}{
      \Axiom$!A \fCenter !A$
      \RightLabel{\RightR{\lnot}}
      \UnaryInf$\fCenter !A, \lnot !A$
      \RightLabel{$\Box*$}
      \UnaryInf$\fCenter \Box!A, \Box\lnot !A$
      \doubleLine
      \RightLabel{\RightR{\lor}}
      \UnaryInf$\fCenter \Box!A \lor \Box \lnot !A$
      \DisplayProof}{}
    \iftag{notprvBox,notprvDiamond}{}{\qquad}
    \iftag{prvDiamond}{
      \Axiom$!A \fCenter !A$
      \RightLabel{\LeftR{\lnot}}
      \UnaryInf$\lnot !A, !A \fCenter$
      \RightLabel{$\Diamond*$}
      \UnaryInf$\Diamond\lnot !A,\Diamond!A \fCenter $
      \RightLabel{\RightR{\lnot}}
      \UnaryInf$\Diamond!A \fCenter \lnot\Diamond\lnot !A$
      \RightLabel{\RightR{\lif}}
      \UnaryInf$ \fCenter \Diamond!A \lif \lnot\Diamond\lnot !A$
      \DisplayProof}{}
  \]
అయితే \iftag{prvBox}{$\Box!A \lor \Box \lnot !A$\iftag{prvDiamond}{ మరియు
$\Diamond!A \lif \lnot\Diamond\lnot !A$ రెండూ}{ఇది}}{$\Diamond!A \lif
\lnot\Diamond\lnot !A$ ఇది} $\Log{K}$లో చెల్లుబాటు కాదు.

పూర్వపక్షంలో~$!A$తో పాటు ఇతర సూత్రాలనూ అనుమతించి,
\iftag{prvBox}{$\Box$ నియమం కుడివైపున~$!A$కు
మాత్రమే $\Box$ చేర్చడానికి అనుమతిస్తే\iftag{prvDiamond}{,
లేదా }{}}{}\iftag{prvDiamond}{$\Diamond$ నియమం
ఎడమవైపున~$!A$కు మాత్రమే $\Diamond$ చేర్చడానికి
అనుమతిస్తే}{} (ఇతర సూత్రాలను మాత్రం మార్చకుండా)
కింది వాటిని \tetoken{వ్యుత్పాదించగలం}{!!{derive}}:
\[
  \iftag{prvBox}{
    \Axiom$!A \fCenter !A$
    \RightLabel{\RightR{\lnot}}
    \UnaryInf$\fCenter !A, \lnot !A$
    \RightLabel{\RightR{\Exchange}}
    \UnaryInf$\fCenter \lnot !A, !A$
    \RightLabel{$\Box*$}
    \UnaryInf$\fCenter \lnot!A, \Box !A$
    \doubleLine
    \RightLabel{\RightR{\lor}}
    \UnaryInf$\fCenter \lnot!A \lor \Box !A$
    \DisplayProof}{}
  \iftag{notprvBox,notprvDiamond}{}{\qquad}
  \iftag{prvDiamond}{
    \Axiom$!A \fCenter !A$
    \RightLabel{\LeftR{\lnot}}
    \UnaryInf$\lnot !A, !A \fCenter$
    \RightLabel{$\Diamond*$}
    \UnaryInf$\Diamond\lnot !A, !A \fCenter $
    \RightLabel{\RightR{\lnot}}
    \UnaryInf$!A \fCenter \lnot\Diamond\lnot !A$
    \RightLabel{\RightR{\lif}}
    \UnaryInf$ \fCenter !A \lif \lnot\Diamond\lnot !A$
    \DisplayProof}{}
\]
అయితే \iftag{prvBox}{$\lnot!A \lor \Box !A$ (ఇది $!A
\lif \Box!A$కు తుల్యం)\iftag{prvDiamond}{ మరియు $!A \lif \lnot\Diamond\lnot !A$
రెండూ}{ఇది}}{$!A \lif \lnot\Diamond\lnot !A$ ఇది} $\Log{K}$లో చెల్లుబాటు కాదు.

\end{document}
